% Self-contained proof fragment. Requires amsmath, amssymb, amsthm and
% theorem environments lemma/proposition. No references to older manuscripts.
\section{A finite family of entropy certificates}\label{sec:sparse-dual}
The following bound replaces a discretization of the full feature space by
sparse linear tests. Its logarithmic dependence on the number of observables
is essential when those observables are indexed by all pairs of Boolean
inputs.

Let $(\Omega,\mathcal A,\mu)$ be any probability space, let
$F=(F_1,\ldots,F_D):\Omega\to[-1,1]^D$ be measurable, and write
\[
 \Lambda(\lambda)=\log\int e^{\lambda\cdot F}\,d\mu,
 \qquad
 \Psi_G(q)=\max_{\|\lambda\|_1\leq G}
       \{\lambda\cdot q-\Lambda(\lambda)\},\qquad G>0.
\]
The maximum exists: its domain is compact and its objective is continuous.
The norm $\|\lambda\|_1$ here is the sum of the absolute values of its
$D$ coordinates.

\begin{lemma}[A bounded exponential tilt with controlled error]
\label{lem:bounded-dual-lift}
For every $q\in\mathbb R^D$, there is a maximizer $\lambda_*$ in the
preceding display such that the probability measure
\[
 \frac{d\nu_*}{d\mu}
     =\exp\{\lambda_*\cdot F-\Lambda(\lambda_*)\}
\]
satisfies the exact identity
\begin{equation}\label{eq:dual-lift-identity}
 \Psi_G(q)
   =D(\nu_*\Vert\mu)
      +G\left\|q-\int F\,d\nu_*\right\|_\infty.
\end{equation}
Consequently, if $\Psi_G(q)\leq K$, then
\[
 D(\nu_*\Vert\mu)\leq K,
 \qquad
 \left\|q-\int F\,d\nu_*\right\|_\infty\leq K/G.
\]
\end{lemma}
\begin{proof}
Boundedness of $F$ permits differentiation under the integral and gives
$\nabla\Lambda(\lambda_*)=\int F\,d\nu_*=:m_*$. The first-order
condition for a differentiable concave function maximized over the convex
$\ell^1$ ball is
\[
 (q-m_*)\cdot(\lambda-\lambda_*)\leq0
       \qquad(\|\lambda\|_1\leq G).
\]
Taking the maximum of the left linear functional over that ball gives
\[
 \lambda_*\cdot(q-m_*)=G\|q-m_*\|_\infty.
\]
Moreover,
$D(\nu_*\Vert\mu)=\lambda_*\cdot m_*-\Lambda(\lambda_*)$.
Adding these two identities proves \eqref{eq:dual-lift-identity}.
Relative entropy is nonnegative, which proves the stated consequences.
The log density of $\nu_*$ relative to $\mu$ is bounded, so its entropy
is finite. No support, density, or nonsingularity assumption on $\mu$
was used.
\end{proof}

For an integer $s\geq1$, introduce the finite set
\[
 \mathcal V_{G,s}
 =\left\{\frac1s\sum_{j=1}^s a_j:
        a_j\in\{0,Ge_1,-Ge_1,\ldots,Ge_D,-Ge_D\}\right\}.
\]
Thus $\#\mathcal V_{G,s}\leq(2D+1)^s$ and every vector in this set has
$\ell^1$ norm at most $G$.

\begin{lemma}[Sparse approximation of the dual value]
\label{lem:sparse-dual-net}
For every $q\in\mathbb R^D$,
\begin{equation}\label{eq:sparse-dual-net}
 \Psi_G(q)\leq\frac{G^2}{2s}
   +\max_{v\in\mathcal V_{G,s}}
        \{v\cdot q-\Lambda(v)\}.
\end{equation}
\end{lemma}
\begin{proof}
Fix $\lambda$ with $\|\lambda\|_1\leq G$. Let the random vector $A$
take the value $G\operatorname{sign}(\lambda_i)e_i$ with probability
$|\lambda_i|/G$ for each nonzero coordinate, and take value zero with the
remaining probability. Then $\mathbb EA=\lambda$.
Take independent copies $A_1,\ldots,A_s$ and put
$V=s^{-1}\sum_j A_j\in\mathcal V_{G,s}$.
For each fixed $\omega\in\Omega$, the scalar $A\cdot F(\omega)$ lies
in $[-G,G]$. Hoeffding's exponential-moment inequality therefore gives
\[
 \mathbb E\exp\{(V-\lambda)\cdot F(\omega)\}
 \leq\exp\{G^2/(2s)\}.
\]
Tonelli's theorem and Jensen's inequality imply
\begin{align*}
 \mathbb E\Lambda(V)
 &=\mathbb E\log\int e^{V\cdot F}\,d\mu\\
 &\leq\log\mathbb E\int e^{V\cdot F}\,d\mu\\
 &\leq\Lambda(\lambda)+G^2/(2s).
\end{align*}
Since $\mathbb E(V\cdot q)=\lambda\cdot q$, at least one realization
$v\in\mathcal V_{G,s}$ satisfies
\[
 v\cdot q-\Lambda(v)
 \geq\lambda\cdot q-\Lambda(\lambda)-G^2/(2s).
\]
Take $\lambda$ to be a maximizer defining $\Psi_G(q)$.
\end{proof}

\begin{proposition}[A nonasymptotic entropy certificate]
\label{prop:finite-entropy-certificate}
Let $X_1,\ldots,X_N$ be independent with law $\mu$, and put
$q_N=N^{-1}\sum_{i=1}^N F(X_i)$. For every $K>0$, $G>0$, and integer
$s\geq1$,
\begin{equation}\label{eq:finite-entropy-tail}
 \Pr\{\Psi_G(q_N)>K\}
 \leq\exp\left\{s\log(2D+1)
          -N\left(K-\frac{G^2}{2s}\right)\right\}.
\end{equation}
Outside this exceptional event, the empirical moment vector is within
$K/G$ in maximum norm of the moments of an exponential tilt with parameter
$\ell^1$ norm at most $G$ and relative entropy at most $K$.
In particular, choosing $s=\lceil G^2\rceil$ bounds the right side by
$\exp\{\lceil G^2\rceil\log(2D+1)-N(K-1/2)\}$.
\end{proposition}
\begin{proof}
For each fixed $v\in\mathcal V_{G,s}$, independence gives exactly
\[
 \mathbb E\exp\{N[v\cdot q_N-\Lambda(v)]\}=1.
\]
Markov's inequality bounds the probability that its bracket exceeds
$K-G^2/(2s)$ by $\exp\{-N(K-G^2/(2s))\}$.
The sparse approximation lemma and a union bound over
$\mathcal V_{G,s}$ prove \eqref{eq:finite-entropy-tail}.
The bounded-dual lift lemma gives the final assertion.
\end{proof}

\paragraph{Application to hidden-layer covariances.}
Take $m$ input points, $D=m^2$, $\mu=\mathcal N(0,Q)$ for any positive
semidefinite covariance $Q$, and $F_{ij}(z)=\tanh(z_i)\tanh(z_j)$,
indexed by all ordered pairs $(i,j)$. Reshape a maximizing parameter
$\lambda_*$ into a matrix and replace it by its symmetric part $\Gamma$.
This preserves its quadratic form and does not increase its entrywise
$\ell^1$ norm. The lifted covariance is then exactly
\[
 \mathsf F_\Gamma(Q)=
 \frac{\mathbb E_{Z\sim\mathcal N(0,Q)}
    [\tanh(Z)\tanh(Z)^\top e^{\tanh(Z)^\top\Gamma\tanh(Z)}]}
      {\mathbb E_{Z\sim\mathcal N(0,Q)}
    e^{\tanh(Z)^\top\Gamma\tanh(Z)}}.
\]
Thus, except with the probability in \eqref{eq:finite-entropy-tail}, the
empirical activation covariance $\widehat Q$ satisfies
\[
 \|\widehat Q-\mathsf F_\Gamma(Q)\|_{\max}\leq K/G,
 \qquad \|\Gamma\|_{1,\mathrm{ent}}\leq G.
\]
This assertion is valid conditionally on any preceding network layers,
using their empirical covariance as $Q$. Its probability bound is uniform
in that covariance, including singular covariances. On the Boolean cube,
$m=2^d$ and the logarithmic factor is
$\log(2m^2+1)\leq(2d+2)\log2$, which is linear in $d$.
